\section{问题分析}\label{sec:analysis}

\subsection{总体分析}

四个问题形成由现象认识到实验决策的递进关系。第一问关注单个催化剂组合内部的温度
关系及一次实验的时间变化；第二问转向组合与温度两类差异的比较；第三问把转化率和
选择性合成为收率，并据此选择实验条件；第四问再把前三问尚未解决的证据不足转化为
有限次数的补充实验。后问继承前问的数据口径和结论边界，但研究对象和所需证据不同，
不能机械套用同一模型。

\subsection{问题一分析}

第一问的两份附件具有不同的数据结构。附件1记录多个催化剂组合在若干固定温度下的
观测，任务是分别说明各组合内部的温度关系；附件2则记录同一次实验随时间的变化，
不能把多个时刻当作相互独立的重复实验。因此，两部分必须分别建立与其观测单位相符的
关系描述。

\subsubsection{附件1：各催化剂组合的温度关系}

每种组合的温度点较少，复杂曲线容易把个别波动解释成稳定规律；不同组合的配方和温度
覆盖又不完全相同，也不宜合并成一条总体曲线。为同时保留局部变化和总体方向，本文先
比较同一组合的相邻温度观测，再用简单线性关系概括其平均变化。该方法只描述各组合
已测温度范围内的经验关系，不承担连续温度预测或因果解释。

\subsubsection{附件2：一次实验的时间变化}

附件2的观测时刻间隔不等，变化量必须结合实际时间间隔解释。与此同时，各类产物
选择性共同构成产物份额，单类选择性的升降可能来自份额重新分配，不能彼此割裂。
因此，本文分别考察主要指标的时间变化和产物份额结构；反应路线文献只用于提出与现象
相容的可能机制，不作为本题具体机理的证明
\cite{aitchison1982,ho2016,moteki2016,sushkevich2017}。

\subsection{问题二分析}

第二问既要判断温度变化，也要比较催化剂组合，还要考察不同组合随温度变化的幅度是否
一致。
这三项任务的比较对象不同：温度影响应在同一组合内部判断，组合差异应在同一温度下
判断，而两者共同变化需要在前两类比较的基础上再作概括。

附件在不同温度下覆盖的组合并不完全相同。若直接比较各温度的全部记录，均值变化会
同时混入温度差异和样本组成变化。为保持比较对象一致，主体分析应使用组合覆盖相同的
温度记录，其他不完整温度只能在具有匹配观测的组合内作补充。同时，题目给出的每个
催化剂组合作为一个完整类别，现有数据不足以把组合差异归因于某个单一配方因素
\cite{searle1980}。

同组配对与同温比较能够回答具体方向和相对高低，但还不能统一衡量组合差异、温度差异
及未被简单相加关系解释的变化。因此，需要在可比数据上建立加和分解。由于每个实验
条件没有重复观测，未解释部分不能单独归因于组合与温度的共同作用或实验误差；模型只作描述性
分解，并通过保持比较结构的删减检查限定结论范围
\cite{nisttwoway,graphpadnorep}。

\subsection{问题三分析}

第三问要求在一般情景和严格低温情景下选择使 C4 烯烃收率尽可能高的条件。催化剂组合
是离散类别，各组合的温度观测又少且不完整。若直接拟合连续曲线再求极值，模型形式会
替代实验数据决定答案；若完全忽略内部缺测，又无法判断已有温度标签中的空白是否可能
改变推荐。

因此，本文把实测条件比较与缺测核查分开：正式推荐只在符合情景约束的真实观测中
产生，内部估计只用于检查缺测单元是否可能改变该推荐。辅助估计不得跨组合借用曲线、
超出组合自身实测范围或生成新配方，也不得以预测值替代实验事实。这样的分层使正式
答案的证据身份明确，同时保留对数据覆盖不足的必要检查。

\subsection{问题四分析}

第四问的目标不是用新增实验重复证明前三问，而是在5次预算内优先减少最影响条件选择的
不确定性。前三问暴露的数据不足可以归为三类：关键温度区间观测稀疏、最高观测缺少
独立重复、不同组合缺少新增同温条件下的直接比较。因此，实验设计既要补充局部温度
信息，也要兼顾可复现性和竞争组合比较，不能只按某个拟合模型的预测收率排序。

要刻画各组合在整个温度范围内的弯曲变化，即建立完整响应面，需要更多实验；本题的
催化剂组合又是不可随意拆分的离散类别。5次预算不足以同时确定全温区的曲线形态、
组合差异和随机误差。本文因而采用证据缺口驱动的局部增补：
先由前三问已经得到的结果明确每项未决问题，再为每次实验规定明确且可核查的职责。
具体缺口、实验条件和观测后的更新规则均在第四问模型建立与求解中给出
\cite{nistreplication,nistrsmcompare}。

